\documentclass[tikz,border=6pt]{standalone}
\usepackage{ctex}
\usepackage{amsmath}
\usepackage{amssymb}
\usepackage{pgfplots}
\pgfplotsset{compat=1.18}
\usetikzlibrary{calc,angles,quotes,arrows.meta,positioning,decorations.markings}
\begin{document}
\begin{tikzpicture}
\begin{axis}[
  width=13cm, height=9cm,
  xmin=0, xmax=1.12, ymin=-0.9, ymax=7.2,
  axis lines=middle,
  xlabel={$p$（事件概率）}, ylabel={$I=-\log_2 p$（bit）},
  xlabel style={at={(axis description cs:0.5,-0.02)},anchor=north},
  ylabel style={at={(axis description cs:0.04,1)},anchor=north west},
  xtick={0.125,0.25,0.5,1}, xticklabels={$\tfrac18$,$\tfrac14$,$\tfrac12$,$1$},
  ytick={1,2,3,4,5,6,7},
  tick label style={font=\small},
  domain=0.0085:1, samples=400,
  clip=false,
]
  % self-information curve I(p) = -log2 p
  \addplot[very thick, blue] {-ln(x)/ln(2)};

  % asymptote p -> 0
  \draw[gray, dashed] (axis cs:0,0) -- (axis cs:0,7.2);
  \node[gray, font=\small, anchor=west] at (axis cs:0.015,6.6)
    {$p\to0^+\Rightarrow I\to+\infty$};

  % point (1,0): certain event, zero information
  \addplot[only marks, mark=*, mark size=2.4pt, black] coordinates {(1,0)};
  \node[anchor=south west, font=\small] at (axis cs:0.78,0.18) {$(1,\,0)$};
  \node[anchor=west, font=\footnotesize, align=left] at (axis cs:0.80,1.05)
    {确定事件\\信息量为 $0$};
  \draw[gray] (axis cs:0.93,0.82) -- (axis cs:0.995,0.10);

  % sample dots at p = 1/8, 1/4, 1/2 with their bit values
  \addplot[only marks, mark=*, mark size=2pt, red!80!black]
    coordinates {(0.5,1) (0.25,2) (0.125,3)};
  \node[red!80!black, font=\footnotesize, anchor=west] at (axis cs:0.5,1.0) {$\,1$ bit};
  \node[red!80!black, font=\footnotesize, anchor=west] at (axis cs:0.25,2.05) {$\,2$ bit};
  \node[red!80!black, font=\footnotesize, anchor=west] at (axis cs:0.135,3.05) {$\,3$ bit};

  % annotation box (semi-transparent), placed in empty middle-right
  \node[draw, rounded corners, fill=yellow!18, fill opacity=0.80, text opacity=1,
        align=left, font=\small, inner sep=5pt]
    at (axis cs:0.74,4.4)
    {越罕见越"惊讶"：\\ 概率 $p$ 越小，\\ 自信息 $I=-\log_2 p$ 越大};
\end{axis}
\end{tikzpicture}
\end{document}
